% !TeX program = pdflatex
% Draf seminar proposal. Isi identitas dan persetujuan resmi sebelum pengajuan.
\documentclass[12pt,a4paper,onecolumn,oneside,final]{report}
\input{config/if-itb-thesis.sty}
\ifproposalbiblatex
  \addbibresource{proposal-references.bib}
\fi

\title{Optimasi Komunikasi Federated Learning Asinkron untuk Network Intrusion Detection melalui Demand-Aware Adaptive Resource Allocation pada Data Non-IID}
\author{Muhammad Fadil Alfaizi}
\nim{123450115}
\date{2026}

\begin{document}
\sloppy
\setlength{\headheight}{15pt}
\pagenumbering{gobble}
\begin{titlepage}
  \centering
  \includegraphics[width=2.1cm]{figures/itera-logo.png}\par
  \vspace{1.8cm}
  {\large\bfseries \MakeUppercase{\thetitle}\par}
  \vfill
  {\large\bfseries LAPORAN PROPOSAL TUGAS AKHIR\par}
  \vspace{0.8cm}
  Disusun untuk Seminar Proposal Tugas Akhir\par
  \vfill
  \theauthor\par
  NIM \printnim\par
  \vfill
  {\bfseries Program Studi Sains Data\par
  Fakultas Sains\par
  Institut Teknologi Sumatera\par
  Lampung Selatan\par}
  \vspace{0.8cm}
  \thedate\par
\end{titlepage}

\pagenumbering{roman}
\tableofcontents
\clearpage
\pagenumbering{arabic}
\renewcommand{\chaptername}{BAB}
\renewcommand{\tablename}{Tabel}
\renewcommand{\figurename}{Gambar}
\renewcommand{\thechapter}{\Roman{chapter}}
\renewcommand{\thesection}{\arabic{chapter}.\arabic{section}}
\justifying

\chapter{Pendahuluan}
\label{bab:pendahuluan}

\section{Latar Belakang}
Sistem deteksi intrusi jaringan (\textit{network intrusion detection system}, NIDS) mempelajari pola lalu lintas untuk membedakan aktivitas normal dan serangan. Data jaringan dapat tersebar di beberapa lokasi dengan distribusi dan kapasitas yang berbeda. \textit{Federated learning} (FL) memungkinkan klien melatih model secara lokal dan mengirim informasi model untuk membentuk model bersama, tanpa memindahkan seluruh data latih mentah ke server \cite{mcmahan2017,agrawal2021}. Akan tetapi, pertukaran model yang berulang menimbulkan biaya komunikasi. Dalam penggunaan FL untuk deteksi intrusi, efisiensi komunikasi harus dinilai bersama kualitas deteksi, terutama kemampuan mengenali serangan.

Pada FL asinkron, server dapat menerapkan pembaruan yang tiba tanpa menunggu seluruh klien. Cara ini mengurangi keharusan sinkronisasi, tetapi pembaruan dari klien yang lebih lambat dapat dihitung dari versi model yang sudah tertinggal \cite{xie2019}. Selain itu, klien yang cepat dapat lebih sering berkontribusi, sedangkan klien lain jarang memperoleh kesempatan. Kondisi tersebut penting ketika jumlah data, waktu pelatihan, dan keadaan jaringan setiap klien tidak sama. Seleksi klien berdasarkan sumber daya sudah dikaji, misalnya oleh Nishio dan Yonetani \cite{nishio2019}, sementara pengurangan ukuran pembaruan melalui sparsifikasi juga telah dipelajari \cite{lin2018}. Kedua keputusan itu memengaruhi biaya komunikasi melalui jalur yang berbeda: seberapa sering klien bertukar pesan dan seberapa besar pesan yang dikirim.

Penelitian terkait deteksi anomali jaringan oleh Marfo dkk. memadukan seleksi klien adaptif, ukuran \textit{batch}, dan pembaruan asinkron; biaya komunikasi yang mereka laporkan terutama berupa waktu \cite{marfo2025}. Penelitian ini mengambil ruang lingkup yang lebih sempit dan terkontrol: satu estimasi kebutuhan komunikasi dari keadaan klien dipakai untuk mengatur \emph{kesempatan berpartisipasi} dan \emph{anggaran ukuran pembaruan}. Efek masing-masing keputusan dan efek gabungannya akan dipisahkan melalui empat kebijakan pada fondasi FL asinkron yang sama. Penghematan akan dihitung dari jumlah \textit{byte} pesan yang benar-benar diserialisasi, termasuk indeks dan metadata, sehingga proporsi parameter yang dipilih tidak langsung dianggap setara dengan penghematan komunikasi.

Data utama yang direncanakan adalah UNSW-NB15 dengan tugas klasifikasi biner normal atau serangan. Dataset ini menyediakan pembagian latih dan uji yang diterbitkan oleh UNSW \cite{moustafa2015,unswdata}. Data latih dibagi kepada sepuluh klien logis menggunakan kemencengan label Dirichlet dengan $\alpha=0{,}5$. Pembagian itu mensimulasikan data \textit{non-IID}, bukan merepresentasikan distribusi jaringan organisasi nyata. Dengan pengaturan tersebut, penelitian menilai apakah keputusan komunikasi yang berubah mengikuti keadaan klien dapat mengurangi biaya tanpa mengorbankan kualitas deteksi di luar batas yang telah ditentukan sebelum eksperimen utama.

\section{Rumusan Masalah}
Permasalahan utama penelitian ini adalah belum diketahuinya manfaat pengendalian bersama partisipasi dan ukuran pembaruan berdasarkan keadaan klien pada FL asinkron untuk NIDS, ketika biaya dihitung pada tingkat \textit{byte} dan kualitas deteksi dijaga. Pertanyaan penelitian dirumuskan sebagai berikut.
\begin{enumerate}
  \item Bagaimana keadaan klien yang tersedia pada saat keputusan dibuat dipetakan menjadi kebutuhan partisipasi dan kebutuhan anggaran pembaruan yang dapat diaudit?
  \item Apakah kebijakan gabungan dapat mengurangi \textit{byte} komunikasi dan waktu untuk mencapai kualitas validasi yang ditetapkan, dibandingkan kebijakan asinkron seragam, tanpa menurunkan kualitas deteksi melebihi margin yang disepakati?
  \item Berapa kontribusi pengaturan partisipasi dan pengaturan anggaran pembaruan secara terpisah terhadap perubahan biaya komunikasi, waktu, dan kualitas model?
  \item Bagaimana heterogenitas jaringan dan pembaruan yang tertinggal berkaitan dengan keputusan kebijakan, pemerataan partisipasi, dan kegagalan mencapai target?
\end{enumerate}

\section{Tujuan Penelitian}
\begin{enumerate}
  \item Merancang pemetaan keadaan klien menjadi dua skor kebutuhan dan dua tindakan komunikasi yang deterministik serta dapat direproduksi.
  \item Mengukur perbedaan \textit{byte} komunikasi, waktu, staleness, partisipasi, dan kualitas deteksi antara kebijakan dasar dan kebijakan gabungan pada data \textit{non-IID} yang tersimulasikan.
  \item Mengisolasi pengaruh setiap tindakan melalui perbandingan empat kebijakan dengan model, agregasi, pelatihan lokal, partisi data, dan profil jaringan yang setara.
\end{enumerate}

\section{Batasan Penelitian}
Penelitian berfokus pada satu model MLP untuk klasifikasi biner UNSW-NB15. Sepuluh klien adalah partisi logis dari dataset, bukan sepuluh institusi pemilik data. Kondisi utama memakai $\alpha=0{,}5$ dan tiga \textit{seed} yang sama pada setiap kebijakan. Model, optimisasi lokal, batas komputasi, serta aturan agregasi dan staleness dipertahankan sama pada seluruh perbandingan. Sumber daya yang diatur adalah kesempatan partisipasi dan ukuran pembaruan \emph{unggah}; jumlah \textit{epoch}, ukuran \textit{batch}, CPU, dan memori tidak menjadi tindakan adaptif. Pengujian dilakukan pada lingkungan penelitian terkendali, sehingga hasilnya tidak langsung menyatakan kinerja pada jaringan produksi. Pengulangan pada dataset lain dan pembanding FL sinkron menjadi perluasan setelah eksperimen utama selesai.

\section{Manfaat dan Luaran}
Manfaat yang diharapkan adalah rancangan pengukuran komunikasi yang dapat diaudit dan bukti empiris tentang kapan pengaturan bersama dua sumber daya komunikasi berguna atau gagal. Luaran Seminar Proposal adalah laporan Bab I--III ini beserta rancangan eksperimennya. Luaran akhir Tugas Akhir direncanakan berbentuk \emph{artikel ilmiah} yang menyajikan metode, hasil perbandingan, ablasi, dan batas validitas setelah eksperimen terlaksana. Tidak ada klaim keberhasilan metode usulan pada tahap proposal.

\chapter{Tinjauan Pustaka}
\label{bab:pustaka}

\section{Network Intrusion Detection dan Dataset}
NIDS berbasis pembelajaran mesin mengklasifikasikan catatan lalu lintas jaringan untuk membantu mengenali aktivitas mencurigakan. Pada penelitian ini, keluaran model hanya dua kelas, yaitu normal dan serangan; model tidak memblokir lalu lintas secara langsung. UNSW-NB15 dipilih karena menyediakan fitur lalu lintas, label normal atau serangan, dan pembagian latih serta uji yang dapat dipakai secara konsisten \cite{moustafa2015,unswdata}. Pembagian tersebut penting agar pengukuran akhir tidak tercampur dengan data yang dipakai untuk memilih aturan atau ambang kebijakan.

Ketidakseimbangan kelas dan perbedaan sebaran antara data latih, validasi, dan uji dapat membuat akurasi saja kurang cukup. Karena itu, \textit{Macro-F1} menjadi ukuran utama kualitas, disertai \textit{recall} kelas serangan dan \textit{area under the precision--recall curve} (AUPRC). Nilai \textit{Macro-F1} adalah rerata F1 kedua kelas:
\begin{equation}
  F1_{\mathrm{macro}}=\frac{F1_{\mathrm{normal}}+F1_{\mathrm{attack}}}{2},\qquad
  F1_c=\frac{2P_cR_c}{P_c+R_c}.
\end{equation}
Di sini $P_c$ dan $R_c$ masing-masing adalah presisi dan \textit{recall} untuk kelas $c$; aturan ketika penyebut nol harus konsisten dalam semua kebijakan.

\section{Federated Learning Asinkron dan Data Non-IID}
Dalam FL, klien menghitung pembaruan dari data lokal, lalu server menggabungkan informasi tersebut menjadi model global \cite{mcmahan2017}. Secara konseptual, tujuan pelatihan adalah meminimalkan $F(w)=\sum_{i=1}^{K}(n_i/n)F_i(w)$, dengan $n_i$ jumlah contoh pada klien $i$ dan $n=\sum_i n_i$. Persamaan ini menyatakan tujuan global, bukan aturan agregasi khusus yang diusulkan penelitian ini. Pada fondasi asinkron, server menerapkan pembaruan saat tiba; jika klien mulai dari versi $v_i^{\mathrm{start}}$ dan pembaruan tiba pada versi $v_i^{\mathrm{arrival}}$, staleness dicatat sebagai
\begin{equation}
  \tau_i=v_i^{\mathrm{arrival}}-v_i^{\mathrm{start}}.
\end{equation}
Aturan penanganan $\tau_i$ harus sama untuk seluruh kebijakan agar perbandingan tetap terisolasi \cite{xie2019}.

Data tiap klien sering memiliki proporsi kelas berbeda. Dalam penelitian ini, pembagian Dirichlet dipakai untuk membentuk kemencengan label yang terkendali. Parameter $\alpha=0{,}5$ dipakai sebagai kondisi utama, dengan manifest partisi dan \textit{seed} tetap. Penggunaan Dirichlet sebagai alat simulasi tingkat heterogenitas telah ditunjukkan pada studi FL sebelumnya \cite{hsu2019}. Karena pembagian dibuat dari satu dataset, kesimpulan hanya berlaku untuk kondisi \textit{non-IID} yang disimulasikan dan profil sistem yang diuji.

\section{Partisipasi Klien dan Anggaran Pembaruan}
Pemilihan klien menentukan siapa yang memperoleh kesempatan mengunduh model dan mengirim pembaruan. FedCS menunjukkan bahwa kapasitas komputasi dan komunikasi klien dapat menjadi dasar seleksi \cite{nishio2019}. Pada FL asinkron, kesempatan partisipasi lebih tepat diatur melalui prioritas atau jeda antar kesempatan berdasarkan peristiwa server daripada melalui jumlah klien per \textit{round}. Pengaturan yang terlalu mengutamakan klien cepat dapat mengabaikan kontribusi klien lain; karena itu distribusi kesempatan dan masa tidak aktif perlu dipantau.

Anggaran pembaruan membatasi bagian perubahan parameter yang dikirim. Salah satu operator yang dapat dipakai adalah \textit{magnitude top-k}, yaitu memilih elemen dengan nilai absolut terbesar dan menyertakan posisi elemennya. Penelitian kompresi gradien telah menunjukkan bahwa komunikasi dapat dikurangi melalui sparsifikasi, tetapi operator dan konteksnya berbeda dari evaluasi FL-NIDS yang direncanakan di sini \cite{lin2018}. Dalam penelitian ini, anggaran kandidat adalah 25\%, 50\%, dan 100\% dari elemen perubahan model. Besar pesan sesungguhnya tetap bergantung pada nilai, indeks, header, dan metadata. Karena itu, penurunan kepadatan elemen tidak diperlakukan sebagai penghematan \textit{byte} sebelum pesan diserialisasi dan dihitung.

\section{Penelitian Terkait dan Posisi Penelitian}
\begin{table}[ht]
\centering
\caption{Ringkasan karya yang melandasi rancangan penelitian}
\label{tab:terkait}
\small
\begin{tabular}{>{\raggedright\arraybackslash}p{2.6cm}>{\raggedright\arraybackslash}p{5.0cm}>{\raggedright\arraybackslash}p{4.9cm}}
\toprule
Karya & Fokus utama & Kaitannya dengan penelitian ini \\
\midrule
McMahan dkk. \cite{mcmahan2017} & Pelatihan model bersama melalui pembaruan lokal. & Dasar FL; bukan uji dua tindakan komunikasi asinkron. \\
Xie dkk. \cite{xie2019} & Optimisasi federasi asinkron dan staleness. & Dasar fondasi asinkron; aturan agregasi dijaga tetap. \\
Nishio dan Yonetani \cite{nishio2019} & Seleksi klien dengan kendala sumber daya. & Landasan tindakan partisipasi; anggaran pembaruan tidak menjadi fokusnya. \\
Lin dkk. \cite{lin2018} & Kompresi gradien untuk mengurangi komunikasi. & Landasan sparsifikasi; biaya pesan NIDS tetap perlu diukur sendiri. \\
Marfo dkk. \cite{marfo2025} & Seleksi adaptif, ukuran \textit{batch}, dan pembaruan asinkron untuk deteksi anomali jaringan. & Karya terdekat; penelitian ini menguji dua tindakan komunikasi dengan ablasi dan akuntansi \textit{byte}. \\
\bottomrule
\end{tabular}
\end{table}

Perbedaan yang diuji penelitian ini bersifat operasional: satu modul estimasi kebutuhan menghasilkan dua keluaran untuk mengatur kesempatan partisipasi dan anggaran pembaruan pada fondasi yang sama. Penelitian ini tidak mengklaim menemukan FL asinkron, seleksi klien, atau sparsifikasi baru. Kesimpulan tentang manfaat gabungan hanya dapat ditarik setelah kebijakan yang mengubah satu tindakan saja dibandingkan dengan kebijakan dasar dan kebijakan gabungan.

\section{Kerangka Pikir dan Hipotesis Kerja}
Keadaan klien yang tersedia pada waktu keputusan membentuk vektor
\begin{equation}
  x_i(t)=\left[n_i,\,T_i^{\mathrm{prev}},\,RTT_i,\,BW_i,\,\tau_i,\,age_i\right].
\end{equation}
Jumlah sampel dan durasi pelatihan sebelumnya menggambarkan beban lokal; \textit{round-trip time} ($RTT_i$) dan \textit{throughput} ($BW_i$) menggambarkan kondisi komunikasi; $\tau_i$ dan $age_i$ menggambarkan ketertinggalan pembaruan dan lama sejak kontribusi terakhir. Satu pemetaan deterministik yang akan ditetapkan setelah kalibrasi menghasilkan $D_i(t)=[d_{\mathrm{part},i},d_{\mathrm{budget},i}]$. Kedua skor dipetakan menjadi prioritas partisipasi dan anggaran pembaruan. Hanya nilai yang telah tersedia sebelum keputusan boleh dipakai, sehingga informasi dari pembaruan atau hasil evaluasi masa depan tidak bocor ke tindakan saat ini.

Hipotesis kerja penelitian adalah kebijakan gabungan dapat mencapai kualitas validasi yang ditetapkan dengan lebih sedikit \textit{byte} daripada kebijakan seragam, sambil mempertahankan \textit{Macro-F1} dalam margin yang ditentukan. Hipotesis ini dapat ditolak. Pengaturan partisipasi saja dan pengaturan anggaran saja diperlukan untuk menjelaskan sumber efek; keluaran negatif tetap dicatat sebagai hasil penelitian.

\chapter{Metode Penelitian}
\label{bab:metode}

\section{Jenis dan Rancangan Penelitian}
Penelitian menggunakan eksperimen kuantitatif komparatif pada prototipe FL asinkron. Variabel perlakuan adalah aturan partisipasi klien dan aturan anggaran pembaruan. Variabel respons mencakup biaya komunikasi, waktu, staleness, pemerataan partisipasi, dan kualitas deteksi. Empat kebijakan membentuk ablasi dua tindakan pada kondisi utama UNSW-NB15, sepuluh klien, kemencengan label Dirichlet $\alpha=0{,}5$, dan sekurang-kurangnya tiga \textit{seed} berpasangan (42, 123, 2026).

Fondasi B0 pada repositori telah dipakai untuk pilot, pencatatan \textit{byte} pesan aplikasi yang diserialisasi, dan pemantauan keadaan dari log peristiwa. Tahap tersebut membuktikan kesiapan sebagian perangkat eksperimen, bukan efektivitas kebijakan gabungan. Pengukur $RTT$ dan \textit{throughput} belum tersedia pada jejak B0 berbasis komunikasi antarproses; keduanya tidak boleh diisi dari waktu pelatihan atau serialisasi. Sebelum eksperimen utama, aturan pengukuran atau simulasi jaringan, normalisasi, perlakuan nilai yang belum tersedia, dan waktu pengambilan keputusan harus dijelaskan serta dibekukan.

\section{Data, Praproses, dan Model}
Penelitian memakai pembagian resmi latih dan uji UNSW-NB15 \cite{unswdata}. Label dipetakan ke normal dan serangan. Data latih dipisah lagi untuk \textit{train-fit} dan validasi; \textit{encoder}, \textit{scaler}, dan transformasi lain hanya dipasang pada \textit{train-fit}. Set uji resmi disimpan untuk evaluasi akhir setelah rancangan kebijakan dan target analisis dibekukan. Partisi sepuluh klien dibuat dari \textit{train-fit} dengan kemencengan label Dirichlet $\alpha=0{,}5$ untuk ketiga \textit{seed}. Manifest partisi memeriksa bahwa setiap baris \textit{train-fit} ditempatkan tepat sekali dan tidak ada baris validasi atau uji yang masuk ke klien.

Satu model \textit{multilayer perceptron} yang sama digunakan oleh semua kebijakan. Konfigurasi kerja menetapkan dua lapisan tersembunyi berukuran 128 dan 64, \textit{dropout} 0,2, presisi FP32, Adam dengan laju belajar 0,001, \textit{batch size} 64, dan satu \textit{local epoch}. Inisialisasi, praproses, optimisasi, batas komputasi, aturan agregasi, dan penanganan staleness harus tetap sama di dalam setiap pasangan perbandingan. Perubahan kebijakan hanya boleh memengaruhi dua tindakan komunikasi yang diteliti.

\section{Kebijakan Pembanding dan Metode Usulan}
\begin{table}[ht]
\centering
\caption{Ablasi kebijakan pada fondasi FL asinkron yang sama}
\label{tab:kebijakan}
\begin{tabular}{>{\raggedright\arraybackslash}p{1.1cm}>{\raggedright\arraybackslash}p{5.7cm}>{\raggedright\arraybackslash}p{5.7cm}}
\toprule
Kode & Kesempatan partisipasi & Anggaran pembaruan unggah \\
\midrule
B0 & Seragam, sesuai aturan dasar & Tetap 100\% \\
B1 & Berdasarkan kebutuhan klien & Tetap 100\% \\
B2 & Seragam, sesuai aturan dasar & Berdasarkan kebutuhan klien: 25\%, 50\%, atau 100\% \\
B3 & Berdasarkan kebutuhan klien & Berdasarkan kebutuhan klien: 25\%, 50\%, atau 100\% \\
\bottomrule
\end{tabular}
\end{table}

Modul usulan membaca keadaan klien yang tersedia sebelum kesempatan berikutnya. Fitur yang dapat dipakai meliputi jumlah sampel, durasi pelatihan terakhir, $RTT$, \textit{throughput}, staleness terakhir, dan lama sejak kontribusi yang diterima. Nilai yang belum pernah diamati tetap diberi status belum tersedia; $RTT$ atau \textit{throughput} tidak diperkirakan dari besaran yang tidak mengukur transfer jaringan. Data kalibrasi digunakan untuk menetapkan skala fitur, arah pengaruhnya, bobot skor, dan ambang kategori rendah, sedang, atau tinggi. Rumus dan penanganan nilai kosong harus deterministik dan dibekukan sebelum eksperimen utama, bersama catatan keadaan $\rightarrow$ skor $\rightarrow$ tindakan untuk setiap keputusan. Label uji dan metrik uji resmi tidak masuk ke masukan modul.

Tindakan partisipasi direncanakan sebagai prioritas atau jeda kelayakan berdasar peristiwa pembaruan server, dengan batas masa tidak aktif untuk mencegah klien terus terabaikan. Nilai batas akhir dipilih melalui kalibrasi dan pilot, lalu dibekukan. Tindakan anggaran memakai \textit{magnitude top-k} pada perubahan parameter yang diunggah; operator, serialisasi, dan aturan rekonstruksi sama pada B2 dan B3. Unduhan model tetap dihitung dalam biaya komunikasi. Eksperimen pendahuluan perlu memeriksa bahwa pilihan 50\% benar-benar menghasilkan ukuran pesan yang sesuai, sebab indeks elemen dapat menghapus penghematan yang tampak dari kepadatan tensor.

\section{Tahapan Eksperimen}
\begin{enumerate}
  \item \textbf{Verifikasi fondasi:} periksa pembagian data, pelatihan terpusat sebagai pemeriksaan kewajaran, penerapan pembaruan saat tiba, perhitungan staleness, rekonstruksi pesan, dan akuntansi \textit{byte}. Hasil pilot B0 dipakai untuk memeriksa instrumen, bukan untuk menyimpulkan kebijakan B3.
  \item \textbf{Kalibrasi keadaan dan kebijakan:} rekam fitur yang tersedia pada saat keputusan, tetapkan aturan pengukuran jaringan, normalisasi, skor kebutuhan, ambang tindakan, dan batas masa tidak aktif. Pilihan dibuat tanpa melihat set uji resmi.
  \item \textbf{Pilot ablasi:} jalankan B0--B3 pada satu \textit{seed} dan satu profil jaringan. Periksa bahwa B1 hanya mengubah partisipasi, B2 hanya mengubah anggaran, dan B3 memakai keduanya. Periksa pula kesesuaian rekonstruksi, log peristiwa, dan kualitas validasi.
  \item \textbf{Pembekuan protokol:} simpan konfigurasi, \textit{seed}, partisi, versi kode, profil jaringan, tingkat konkurensi, aturan penghentian, interval evaluasi, target \textit{Macro-F1}, dan rencana analisis beserta \textit{hash}. Sesudahnya, eksperimen utama tidak diubah berdasarkan hasil yang terlihat.
  \item \textbf{Eksperimen utama:} jalankan B0--B3 pada ketiga \textit{seed} dengan pasangan partisi, inisialisasi, batas komputasi, dan profil jaringan yang setara. Profil normal menjadi kondisi utama; profil jaringan heterogen digunakan untuk analisis pengaruh kondisi sistem setelah rancangan utama stabil. Urutan kedatangan asinkron dicatat karena dapat berbeda meskipun profil sama.
  \item \textbf{Analisis akhir:} hitung respons per \textit{run}, bandingkan kebijakan secara berpasangan, audit kasus kegagalan, lalu evaluasi set uji resmi sekali sesuai protokol yang dibekukan.
\end{enumerate}

\section{Metrik dan Analisis Data}
Untuk setiap pesan $m$, ukuran komunikasi adalah $|\mathrm{serialize}(m)|$ \textit{byte}. Biaya total ialah jumlah ukuran pesan unggah dan unduh, termasuk nilai, indeks, header, serta metadata yang diserialisasi. Biaya ini adalah \emph{pesan aplikasi}; bila overhead transport ingin diklaim, ia harus diukur terpisah. Ukuran primer adalah total \textit{byte} dan \textit{bytes-to-target}, yaitu biaya kumulatif pertama ketika \textit{Macro-F1} validasi mencapai target $\theta$ yang telah dibekukan:
\begin{equation}
  B_{\theta}=\min\{B(v): F1_{\mathrm{macro,val}}(v)\geq\theta\}.
\end{equation}
\textit{Time-to-target} didefinisikan secara sepadan dari waktu berjalan. Run yang tidak mencapai $\theta$ harus dilaporkan sebagai tidak mencapai target beserta batas pengamatan, bukan diberi nilai nol atau dikeluarkan diam-diam. Ukuran lain meliputi jumlah pembaruan yang diterima, distribusi staleness, bagian kontribusi tiap klien, masa tidak aktif maksimum, \textit{recall} serangan, AUPRC, serta waktu serialisasi dan kompresi.

Perbandingan dilakukan pada unit \textit{run/seed}, bukan menganggap semua pembaruan dalam satu \textit{run} sebagai sampel independen. Selisih B3--B0 dan ablasi B1/B2 dilaporkan per \textit{seed}, dengan rerata dan sebaran. Sasaran praktis awal yang tercatat adalah pengurangan biaya komunikasi 15\%, penurunan \textit{Macro-F1} akhir tidak lebih dari satu poin persentase, dan pemburukan \textit{time-to-target} tidak lebih dari 10\%; definisi operasional dan titik evaluasinya akan dibekukan sebelum eksperimen utama. Untuk pengujian \textit{non-inferiority}, margin \textit{Macro-F1} adalah $\delta=0{,}01$ pada selisih B3--B0; kriteria yang direncanakan ialah batas bawah interval kepercayaan 95\% selisih tersebut lebih besar dari $-0{,}01$. Tiga \textit{seed} memberi ketidakpastian besar, sehingga seluruh selisih mentah harus ditampilkan dan klaim statistik dibatasi oleh daya yang tersedia. Grafik biaya kumulatif terhadap kualitas dan waktu serta audit partisipasi dipakai untuk menjelaskan perbedaan kebijakan.

\section{Validitas dan Keterulangan}
Risiko utama meliputi kebocoran informasi uji ke kalibrasi, ketidaksamaan kondisi antar kebijakan, simulasi \textit{non-IID} yang tidak mewakili institusi nyata, dan perbedaan antara \textit{byte} pesan aplikasi dengan penggunaan jaringan sesungguhnya. Setiap run menyimpan konfigurasi, manifest data dan partisi, versi perangkat lunak, log peristiwa, ukuran pesan, keluaran validasi, dan \textit{hash} artefak. Klaim penghematan hanya dibuat untuk biaya yang benar-benar terukur; laporan akhir harus menyebutkan secara eksplisit jika profil jaringan atau overhead transport belum berhasil diukur.


\clearpage
\renewcommand{\bibname}{Daftar Pustaka}
\phantomsection
\addcontentsline{toc}{chapter}{Daftar Pustaka}
\ifproposalbiblatex
  \printbibliography
\else
  \bibliographystyle{ieeetr}
  \bibliography{proposal-references}
\fi
\end{document}
